\documentclass[a4paper,landscape]{article}

\usepackage[utf8]{inputenc}
\usepackage[T1]{fontenc}
\usepackage[italian]{babel}
\usepackage{tikz}
\usetikzlibrary{
    positioning,
    arrows.meta,
    shapes.multipart
}

\pagestyle{empty}

\begin{document}

\begin{center}
\Large\textbf{Diagramma delle classi Ghisa App}
\end{center}

\vspace{0.5cm}

\begin{center}
\begin{tikzpicture}[
    scale=0.82,
    transform shape,
    classe/.style={
        draw,
        rectangle split,
        rectangle split parts=3,
        rectangle split part align={center,left,left},
        text width=3.6cm,
        font=\small,
        inner sep=5pt
    },
    relazione/.style={draw,thick},
    eredita/.style={
        draw,thick,
        -{Triangle[open,length=9pt,width=8pt]}
    }
]

% CLASSI

\node[classe] (utente) at (0,0) {
    \textbf{Utente}
    \nodepart{second}
    - idUtente : int
    \nodepart{third}
    + visualizzaEsercizi() : void\\
    + filtraEsercizi() : void
};

\node[classe] (registrato) at (0,-5) {
    \textbf{UtenteRegistrato}
    \nodepart{second}
    - nome : string\\
    - email : string\\
    - password : string
    \nodepart{third}
    + registrati() : void\\
    + creaScheda() : void\\
    + monitoraCarichi() : void
};

\node[classe] (admin) at (-6,-5) {
    \textbf{Admin}
    \nodepart{second}
    - idAdmin : int
    \nodepart{third}
    + caricaEsercizio() : void\\
    + gestisciCarichi() : void\\
    + gestisciSchede() : void
};

\node[classe] (esercizio) at (7,0) {
    \textbf{Esercizio}
    \nodepart{second}
    - idEsercizio : int\\
    - nome : string\\
    - descrizione : string\\
    - gruppoMuscolare : string
    \nodepart{third}
    + visualizza() : void\\
    + modifica() : void
};

\node[classe] (scheda) at (7,-5) {
    \textbf{Scheda}
    \nodepart{second}
    - idScheda : int\\
    - nome : string\\
    - dataCreazione : date
    \nodepart{third}
    + aggiungiEsercizio() : void\\
    + rimuoviEsercizio() : void\\
    + modificaScheda() : void
};

\node[classe] (carico) at (7,-10) {
    \textbf{Carico}
    \nodepart{second}
    - idCarico : int\\
    - peso : double\\
    - ripetizioni : int\\
    - data : date
    \nodepart{third}
    + registraCarico() : void\\
    + aggiornaCarico() : void\\
    + visualizzaStorico() : void
};

% EREDITARIETA'

\draw[eredita]
(registrato.north) -- (utente.south);

\draw[eredita]
(admin.north) |- (-3,0) -- (utente.west);

% ASSOCIAZIONI

% Utente visualizza esercizi
\draw[relazione]
(utente.east) -- (esercizio.west)
node[midway,above] {visualizza}
node[pos=0.12,below] {0..*}
node[pos=0.88,below] {0..*};

% Utente registrato crea schede
\draw[relazione]
(registrato.east) -- (scheda.west)
node[midway,above] {crea}
node[pos=0.12,below] {1}
node[pos=0.88,below] {0..*};

% Scheda contiene esercizi
\draw[relazione]
(esercizio.south) -- (scheda.north)
node[midway,right] {contiene}
node[pos=0.15,left] {0..*}
node[pos=0.85,left] {1..*};

% Scheda registra carichi
\draw[relazione]
(scheda.south) -- (carico.north)
node[midway,right] {registra}
node[pos=0.15,left] {1}
node[pos=0.85,left] {0..*};

% Utente registrato monitora carichi
\draw[relazione]
(registrato.south) |- (carico.west)
node[pos=0.25,left] {monitora};

\end{tikzpicture}
\end{center}

\end{document}